\subsection{有限长度模的对偶}

设 A 为 Noether 环；用 \(\Omega\) 表示它的极大理想组成的集合。推广上一节给出的定义，我们将说一个 A-模 J 是 Matlis A-模，如果它是内射的，如果它的伴随素理想都是 A 的极大理想，并且如果对 A 的每个极大理想 \(\mathfrak{m}\) ，A/ \(\mathfrak{m}\)-向量空间 \(\mathrm{Hom}_{\mathrm{A}}(\mathrm{A}/\mathfrak{m},\mathrm{J})\) 的维数为 1。对每个 \(\mathfrak{m} \in \Omega\) ，选取 \(\kappa(\mathfrak{m}) \to \mathrm{I}(\mathfrak{m})\) 作为 A-模 \(\kappa(\mathfrak{m})\) 的一个内射包；A-模 \(\bigoplus_{\mathfrak{m} \in \Omega} \mathrm{I}(\mathfrak{m})\) 是一个 Matlis 模，并且每个 Matlis A-模都同构于它（\(n^{\circ}\) 1，定理 1）。

回忆（VIII, § 1, 第 5 节），我们用 \(Z_0(A)\) 表示 \(Z\)-模 \(Z^{(\Omega)}\)，并用 \(\varepsilon : Z_0(A) \to Z\) 表示把基 \(\Omega\) 的每个元素映到 1 的线性型。若 M 是一个有限长度的 A-模，则 \(A_m\)-模 \(M_m\) 对每个 \(m \in \Omega\) 都是有限长度的，且除有限多个理想 \(m \in \Omega\) 外为零。我们令

\[
z _ {0} (M) = \sum_ {\mathfrak {m} \in \Omega} \operatorname{long} _ {\mathrm{A} _ {\mathfrak {m}}} (\mathrm{M} _ {\mathfrak {m}}) [ \mathfrak {m} ] \quad \text { in } \quad Z _ {0} (\mathrm{A});
\]

我们有 \(\mathrm{long}_{\mathrm{A}}(\mathrm{M}) = \varepsilon(z_0(\mathrm{M}))\)（同处，例 3）。反之，一个使 \(\mathrm{long}_{\mathrm{A_m}}(\mathrm{N_m})\) 对每个 \(\mathfrak{m} \in \Omega\) 有限、且在 \(\Omega\) 的有限子集 I 之外为零的 A-模 N 是有限长度的：事实上 N 同构于 \(\bigoplus_{\mathfrak{m} \in I} \mathrm{N}_{\mathfrak{m}}\) 的一个子模（II, § 3, 第 3 节，定理 1 的推论 2），并且我们有 \(\mathrm{long}_{\mathrm{A_m}}(\mathrm{N_m}) = \mathrm{long}_\mathrm{A}(\mathrm{N_m})\)，因为每个 \(\mathrm{A_m}\)-单模都同构于 \(\kappa(\mathfrak{m})\) ，从而作为 A-模是单的。

设 J 为一个 Matlis A-模。对每个 A-模 M，我们将用 \(D_{A}(M)\) 或简记作 \(D(M)\) 表示 A-模 \(\operatorname{Hom}_{\mathrm{A}}(\mathrm{M},\mathrm{J})\)。设 \(\alpha_{M}\) 为从 M 到 \(D(D(M))\) 的同态，它由 \(\alpha_{M}(m)(f)=f(m)\) （对 \(m\in M\)，\(f\in D(M)\) ）所定义。

命题 4.—— A-模 M 是有限长度的，必须且只须 D(M) 是有限长度的。此时我们有 \(z_{0}(\mathrm{D}(\mathrm{M})) = z_{0}(\mathrm{M})\)，\(\operatorname{long}_{\mathrm{A}}(\mathrm{M}) = \operatorname{long}_{\mathrm{A}}\mathrm{D}(\mathrm{M})\)，\(\operatorname{Ann}_{\mathrm{A}}(\mathrm{M}) = \operatorname{Ann}_{\mathrm{A}}(\mathrm{D}(\mathrm{M}))\)，并且 A-线性映射 \(\alpha_{M}\) 是双射。

对每个 \(\mathfrak{m} \in \Omega\)，\(A_{\mathfrak{m}}\)-模 \(D(M)_{\mathfrak{m}}\) 等同于 \(\operatorname{Hom}_{A_{\mathfrak{m}}} (M_{\mathfrak{m}}, J_{\mathfrak{m}})\)（II, § 2, 第 7 节，命题 19）；命题的第一个断言由定理 2, c) 以及上面给出的有限长度模的刻画立得。以下设 M 是有限长度的；我们有 \(\operatorname{long}_{A_{\mathfrak{m}}} (D(M)_{\mathfrak{m}}) = \operatorname{long}_{A_{\mathfrak{m}}} (M_{\mathfrak{m}})\) （对每个 \(\mathfrak{m} \in \Omega\) ，同处），由此得 \(z_0(D(M)) = z_0(M)\) 和 \(\operatorname{long}_A(D(M)) = \operatorname{long}_A(M)\) 。此外，映射 \((\alpha_{M})_{\mathfrak{m}}: M_{\mathfrak{m}} \to D(D(M))_{\mathfrak{m}}\) 等同于典范同态 \(\alpha_{M_{\mathfrak{m}}}\) ，后者是双射（同处）；从而 \(\alpha_{M}\) 是双射。

对每个 \(a \in A\)，我们有 \(D(a_{M}) = a_{D(M)}\)，从而 \(\text{Ann}_{A}(M) \subset \text{Ann}_{A}(D(M))\)。把这个结果应用于 A-模 \(D(M)\) ，就得到相反的包含关系，从而得到等式 \(\text{Ann}_{A}(M) = \text{Ann}_{A}(D(M))\) 。

例.—— 设 A 为主理想环，K 为其分式域。A-模 K/A 是一个 Matlis 模：事实上，从 K/A 到 \(\prod_{\mathfrak{m}\in\Omega}\mathrm{K}/\mathrm{A}_{\mathfrak{m}}\) 的典范映射诱导一个从 K/A 到 \(\bigoplus_{\mathfrak{m}\in\Omega}\mathrm{K}/\mathrm{A}_{\mathfrak{m}}\) 上的同构（A, VII, 第 10 页，定理 2）；于是断言由第 1 节例 1 立得。此外，我们在 A, VII, § 4, 第 9 节中已经证明：当环 A 是主理想环时，映射 \(\alpha_{M}\) 对每个有限长度的 A-模 M 都是双射。

